% =====================================================================================================
% Kapitel 3: Bestandsaufnahme und Planung                                       Zielumfang: ca. 3 Seiten
% Zeigt methodisches Vorgehen (fehlt bei den Vorgängern) und begründet die Maßnahmen in Kapitel 4–6.
% =====================================================================================================

\section{Bestandsaufnahme und Planung}

\subsection{Ist-Zustand der Anlage}
% ZIEL:    Dokumentierter Ausgangszustand vor allen Änderungen.
% INHALT:  - Sicherheitsrelais: lose am Prototyp; Ist-Verdrahtung als Tabelle/Stromlaufplan
%            (S11/S12/S21/S22 Not-Halt, Brücke S12–S34, A1–33–41, 13/14 und 23/24 Motoren, Y32 frei)
%          - Scanner: RS232 über CBX100 an Laptop, COM3, 19200 Baud, 8N1, Start-String „LON“
%          - IO-Link-Master: auf der Schleppkette, Port 1 DTI515, Port 2 OY5100
%          - Netzwerk: Adressplan (SPS 192.168.0.20, …), Switch-Belegung
%          - Fotos „vorher“
% QUELLEN: eigene Aufnahme (Datum!), Vorgängerarbeiten, Code T3200
% UMFANG:  ca. 1,5 Seiten
%
% -----------------------------------------------------------------------------------------------------
% MATERIAL AUS DEM INTERVIEW (05.11.2026) – Stichpunkte, vor dem Ausformulieren prüfen
% -----------------------------------------------------------------------------------------------------
% Aufbau Steuerungsgestell (neben dem Prototyp):
%   - Alu-Gestell mit Einschubplatten, 2 Einschubbereiche; Hersteller vermutlich Christiani (prüfen)
%   - oben links: Druckluft-/Wartungseinheit (versorgt den Zylinder am Regalbediengerät)
%   - oben rechts: Siemens HMI (Typ: TODO)
%   - unten links: Hutschiene mit Siemens S7-1500 (CPU-Typ: TODO)
%   - unten rechts: Netzwerk-Switch auf Hutschiene (Typ: TODO)
%   - NEU geplant rechts neben dem Switch: PILZ PNOZ s4
%   - NEU geplant rechts daneben: Datalogic CBX100 (Anschlussbox des Scanners)
%   - zusätzliches Netzwerkkabel am Switch vorbereitet für die Scanner-Anbindung (Konverter)
% Bewegungen des Regalbediengeräts:
%   - horizontale Achse (nach rechts): Motor
%   - vertikale Achse (nach oben): Motor
%   - Einlagern: Zylinder (Druckluft) fährt vor → Gerät senkt leicht ab → Zylinder fährt ein →
%     Holzklotz/Palette liegt im Regalfach
% Motorschaltung:
%   - Motoren werden über kleine Relais geschaltet (weiterhin verbaut)
%   - zusätzlich sind die Sicherheitskontakte des PNOZ s4 in die Motorleitungen (+) verdrahtet
% Netzwerk: IP-Adressen außer SPS (192.168.0.20) noch unbekannt → aus TIA-Netzsicht ablesen
%
% Vorschlag Tabelle für 3.1 (Werte ergänzen):
\begin{table}[H]
  \centering
  \caption{Komponenten des Steuerungsgestells (Ist-Zustand, Stand 11/2026)}
  \label{tab:komponenten_gestell}
  \begin{tabular}{|l|l|l|l|}
    \hline
    \textbf{Position} & \textbf{Komponente} & \textbf{Typ} & \textbf{IP-Adresse} \\ \hline
    oben links   & Druckluft-/Wartungseinheit & TODO & -- \\ \hline
    oben rechts  & HMI (Siemens)              & TODO & TODO \\ \hline
    unten links  & SPS (Siemens S7-1500)      & TODO & 192.168.0.20 \\ \hline
    unten rechts & Netzwerk-Switch            & TODO & TODO \\ \hline
    neu          & Sicherheitsrelais PILZ PNOZ s4 & TODO Bestellnr. & -- \\ \hline
    neu          & Datalogic CBX100           & CBX100 & -- \\ \hline
    neu          & RS232-Ethernet-Konverter   & TODO & TODO \\ \hline
  \end{tabular}
\end{table}

\subsection{Identifizierte Schwachstellen und Anforderungen}
% ZIEL:    Aus dem Ist-Zustand messbare Anforderungen ableiten.
% INHALT:  - Tabelle: Schwachstelle → Anforderung → Kapitel
%            z. B. automatischer Wiederanlauf → überwachter Start (Kap. 4)
%                  Scanner an Laptop/COM-Port → Netzwerkanbindung, Reconnect (Kap. 5)
%                  Masse auf Schleppkette → feste Montage (Kap. 6)
% UMFANG:  ca. 1 Seite

\subsection{Vorgehen}
% INHALT:  - Reihenfolge der Umsetzung und Begründung (Scanner isoliert testbar; Hardware in einem
%            Umbaufenster), Abstimmung/Freigabe durch Betreuer, Prüfprotokolle
% UMFANG:  ca. 0,5 Seiten
